% ==============================================================================
% Related Work
% ==============================================================================
\section{Related Work}
\label{sec:related}

\paragraph{GUI grounding.}
GUI grounding maps a natural-language instruction to an executable screen location, and is harder than general visual grounding because of dense text, small icons, regular layouts, and high-resolution screenshots~\citep{screenspotpro}. Early screenshot-based agents such as META-GUI~\citep{sun2022metagui}, CogAgent~\citep{hong2024cogagent}, SeeClick~\citep{cheng2024seeclick}, Ferret-UI~\citep{you2024ferretui}, and OmniParser~\citep{lu2024omniparser} showed that visual grounding can replace reliance on HTML, DOM, or accessibility trees, and subsequent work scaled model capacity and action traces, including UGround~\citep{gou2024navigating}, OS-Atlas~\citep{Os-atlas}, ShowUI~\citep{lin2025showui}, Aguvis~\citep{xu2024aguvis}, UI-TARS~\citep{qin2025uitars}, GUI-Owl~\citep{ye2025guiowl,xu2026mobile}, and UI-Venus~\citep{gu2025uivenus,team2026ui}, the four backbones we retrofit. Reinforcement-learning variants such as UI-R1~\citep{lu2025ui}, GUI-R1~\citep{luo2025gui}, InfiGUI-R1~\citep{liu2025infigui}, GUI-G1~\citep{zhou2025gui}, and InfiGUI-G1~\citep{liu2026infigui} further refine coordinate prediction with grounding-specific rewards, but all read out the final autoregressive layer, the assumption our analysis (\S3) revisits.

\paragraph{Coordinate-free and region-centric grounding.}
Because coordinate regression is numerically unstable and its point-level supervision is ambiguous, several methods instead read a region-level posterior. Tree-of-Lens~\citep{fan2024readpointed} builds hierarchical layout structures from user-indicated points; TAG~\citep{xu2025attention}, RegionFocus~\citep{luo2025regionfocus}, DiMo-GUI~\citep{wu2025dimogui}, and FocusUI~\citep{ouyang2026focusui} mine intrinsic attention, rescale the viewport, or select position-preserving visual tokens at test time. GUI-Actor~\citep{wu2025gui} is closest in spirit: a coordinate-free, frozen-backbone head reported on ScreenSpot-Pro with Qwen2/Qwen2.5 backbones. It attaches a single $\sim$25M-parameter action head to the \emph{final} layer with no cross-layer fusion; \textsc{ZwerGe-UI} instead distributes ${\approx}51$--$75$M parameters across independent probes on \emph{multiple intermediate} layers plus a ${\approx}200$K-parameter fusion head, motivated by our finding that the most spatially informative layer is backbone-dependent and typically precedes the final layer (\S3.1), so a single last-layer head cannot recover it. We view \textsc{ZwerGe-UI} as a multi-layer, depth-aware generalization of this coordinate-free family. Re-Prefill~\citep{lin2026happens} is concurrent and complementary: it studies the \emph{stage-level} asymmetry between prefill and decoding, showing the prefill stage already fixes candidate elements, and re-prefills with attention guidance to refine them. We instead study the \emph{layer-level} asymmetry \emph{within} a single prefill, showing that spatial localization and coordinate serialization peak at different depths; the two views address different granularities of the same underlying observation that GUI grounding is not solely a final-layer, autoregressive phenomenon.

\paragraph{Where does grounding live inside the network?}
Our diagnostic lenses build on a broader literature examining what individual layers and tokens encode. Logit-lens probing decodes intermediate hidden states with the frozen output head~\citep{nostalgebraist2020logitlens}, and tuned-lens learns a calibrated per-layer affine map to correct for representational drift before decoding~\citep{belrose2023tunedlens}, the technique we adapt as a control for our serialization lens (\S3.2). Layerwise probing has also been used to localize factual and semantic content in language models~\citep{geva2021transformer,meng2022locating}, and analyses of vision transformers show that final-layer tokens can lose spatial fidelity while a small set of high-norm ``register'' tokens absorb global computation instead of encoding local content~\citep{darcet2024registers,lan2024clearclip,chen2025rethinking}. This echoes attention-sink behavior documented in LLMs~\citep{xiao2023streamingllm} and in large multimodal models~\citep{kang2025visualsink}, where high-attention or high-norm tokens serve internal bookkeeping rather than task-relevant regions. Attention-based explanation methods such as Attention Rollout~\citep{abnar2020attentionflow} and relevance propagation~\citep{chefer2021transformer} likewise assume the final layer is the right place to look; our norm census (supplementary material) shows the same depth-dependent magnitude growth these works report, but for GUI spatial posteriors rather than classification tokens, and ties it directly to a measurable grounding-accuracy collapse rather than a purely representational curiosity. Together, this literature motivates treating layer depth as a first-class design axis rather than reading out only the last layer, which is the premise \textsc{ZwerGe-UI} operationalizes for GUI grounding.

\paragraph{Efficient multimodal grounding backbones.}
GUI screenshots are high-resolution, so the visual token budget is a practical constraint on any read-out design. Foundational vision-language architectures~\citep{dosovitskiy2021vit,radford2021clip,bai2023qwenvl,bai2025qwen25vl,qwen3vl} and open-vocabulary detectors~\citep{carion2020detr,kamath2021mdetr,li2022glip,liu2023groundingdino,minderer2022owlvit,kirillov2023sam} establish the patch-token representations our probes read from, while token-compression methods such as ViLT~\citep{kim2021vilt}, PuMer~\citep{cao2023pumer}, FastV~\citep{chen2024fastv}, and TokenPacker~\citep{li2024tokenpacker} reduce the same visual redundancy that our patch-level posterior operates over. Our probes are complementary to this line: rather than compressing tokens before the backbone processes them, we read an existing intermediate representation after a single frozen forward pass, so the retrofit composes with any of these efficiency techniques without retraining the backbone.
